\section{Contratos, gobernanza y reglas económicas}
\subsection{Contratos del protocolo activo}
\begin{longtable}{L{.28\linewidth}L{.62\linewidth}}
\caption{Responsabilidades de los cinco contratos activos.}\label{tab:contratos}\\
\toprule Contrato & Responsabilidad \\ \midrule\endfirsthead
\toprule Contrato & Responsabilidad \\ \midrule\endhead
\code{TimelockController} & Programación y ejecución con retraso y roles de gobernanza.\\
\code{CommunityRegistry} & Ciclo de vida, membresías, nodos, medidores, precio y vínculos con los contratos comunitarios.\\
\code{SUNToken} & ERC-20 con emisión restringida, 18 decimales, pausado, quema y anclajes contables.\\
\code{CommunitySettlement} & Lotes, revisión, finalización, prueba de inyección, emisión acotada y retiro por consumo.\\
\code{NativeSettlementVault} & Depósitos de ETH y SUN, reservas, entrega atómica y retiros por el propietario.\\
\bottomrule
\end{longtable}
Los contratos utilizan Solidity 0.8.24 y componentes OpenZeppelin. Los ABI nativos compartidos entre backend y frontend se sincronizan mediante un script específico. La administración se asigna al contrato timelock. El minter de SUN se concede a CommunitySettlement y se revoca al otorgante inicial una vez autorizado el protocolo \citep{blockchain,contracts}.

\subsection{Ciclo de vida y gobernanza}
La comunidad pasa por propuesta, revisión, aprobación y activación, con transiciones autorizadas para rechazo, suspensión o archivo. La membresía también tiene estados y no se obtiene por conectar una wallet. El operador comunitario y los roles de registrador de medidores, oráculo, liquidación y pausa son permisos diferentes.

Una acción del timelock se \textbf{programa} y después se \textbf{ejecuta} al cumplirse el retraso. El entorno local utiliza demora cero y conserva las dos transacciones de programación y ejecución. El lote también tiene un periodo de revisión antes de su finalización.

El precio por kWh se establece mientras la comunidad está aprobada y queda fijado al activarla. Ese precio se incorpora a las asignaciones del lote.

\subsection{Asignación de cantidades}
Sea $S$ la inyección neta elegible total y $D$ la demanda neta elegible. La cantidad asignable cumple:
\begin{equation}
A=\min(S,D),\quad 0\leq A\leq S,\quad A\leq D.
\end{equation}
La política compensa importación y exportación del intervalo y la autooferta de varios nodos de una misma wallet. La política permite aplicar un factor de pérdidas a la oferta elegible.

Para repartir $A$ entre demandas $d_j$, la base entera es:
\begin{equation}
a_j^{(0)}=\left\lfloor\frac{A d_j}{\sum_k d_k}\right\rfloor.
\end{equation}
Las unidades restantes se distribuyen por mayores residuos y desempate estable de identidad. Se realiza un reparto análogo de oferta y un emparejamiento determinista productor-consumidor. El reparto se realiza entre participantes distintos. El uso de Wh enteros evita residuales energéticos de punto flotante \citep{allocation,settlement}.

\subsection{Lotes y compromisos de evidencia}
El oráculo presenta intervalo, política, raíz Merkle, totales y asignación en un lote. El estado inicial en cadena es \code{CHALLENGEABLE}. Después de la revisión, la gobernanza puede finalizarlo o rechazarlo. Solo un lote canónico puede quedar finalizado por comunidad e intervalo bajo las reglas actuales.

Una prueba Merkle demuestra pertenencia a la raíz comprometida. Las hojas incluyen clase de operación, lote, claim o asignación, nodos, medidores, beneficiarios y cantidades. La doble aplicación de keccak y el orden de pares coinciden entre ABI, backend y contratos. La raíz permite comprobar que una prueba corresponde al conjunto aprobado.

El contrato permite sustituir un lote finalizado que conserve las condiciones de uso requeridas.

\subsection{Emisión, unidades y precio}
La equivalencia contable es:
\begin{equation}
1\;\text{SUN}=1\;\text{kWh},\qquad
1\;\text{Wh}=10^{15}\;\text{unidades base SUN}.
\end{equation}
La finalización del lote habilita la emisión. El trabajador tramita la transacción y consulta su confirmación. La emisión agregada no puede exceder la inyección aprobada y cada claim es único. La emisión precede a la financiación y liquidación de las asignaciones.

El precio $p$ procede de la política comunitaria. Para una asignación de $w$ Wh, el coste en wei es:
\begin{equation}
c=\left\lceil\frac{w\,p}{1000}\right\rceil.
\label{eq:coste}
\end{equation}
En el ejemplo, $p=10^{15}$ wei/kWh, equivalente a 0,001 ETH/kWh. Así, 108 Wh corresponden a 0,108 SUN y 0,000108 ETH. Fraccionar cantidades en varias parejas puede añadir hasta un wei de redondeo por pareja frente a una única cuenta conjunta; el mismo importe se acredita y se adeuda, sin residuo monetario sin propietario.

\subsection{Depósitos, reserva y entrega}
El productor primero autoriza gasto ERC-20 al vault y después deposita SUN. El consumidor deposita ETH nativo y establece un límite de precio. El participante firma el aporte de ETH desde su wallet. Una reserva exige prueba de asignación, cantidad, precio, ETH y SUN suficientes; si falta algún respaldo, la transacción revierte completa.

La ejecución de una reserva válida entrega SUN al consumidor y acredita ingresos retirables del productor en la misma transacción. El ingreso se acredita durante la liquidación y se cobra con \code{withdrawEarnings}. La cancelación de una reserva libera los activos correspondientes.

\begin{figure}[htbp]
\caption{Flujo de activos durante la liquidación}\label{fig:activos}

\centering
\begin{tikzpicture}[node distance=.85cm and 1.2cm,box/.append style={text width=3.6cm}]
\node[box] (prod) {Productor\\SUN en wallet};
\node[box,right=of prod] (v) {Vault\\SUN y ETH reservados};
\node[box,right=of v] (cons) {Consumidor\\ETH en wallet};
\node[box,below=1.2cm of prod] (earn) {Ingresos acreditados\\y retiro de ETH};
\node[box,below=1.2cm of v] (exec) {Ejecución atómica\\de la asignación};
\node[box,below=1.2cm of cons] (burn) {SUN recibidos\\y retiro por consumo};
\draw[arr] (prod)--node[labelnode,above]{autorizar y depositar}(v);
\draw[arr] (cons)--node[labelnode,above]{depositar ETH}(v);
\draw[arr] (v)--(exec);
\draw[arr] (exec)--(earn);
\draw[arr] (exec)--(burn);
\end{tikzpicture}
\par\noindent\small\textit{Nota.} Elaboración propia a partir del protocolo implementado. La quema posterior de SUN registra el consumo y no añade un pago ETH \citep{settlement}.

\end{figure}

\subsection{Retiro por consumo e invariantes}
El consumidor aprueba un allowance diferente, dirigido a CommunitySettlement, y presenta una prueba de consumo. El retiro exige lote finalizado, participante/nodo/medidor autorizados y entrega financiera completa de la asignación comprometida para ese nodo. Se queman los SUN y se marca el consumo como utilizado. La prueba vincula los SUN utilizados con el consumo entregado.

Los invariantes principales de respaldo son:
\begin{align}
\operatorname{ETH}(vault)&\geq \sum_u disponible_u+\sum_u reservado_u+\sum_u ingresos_u,\\
\operatorname{SUN}(vault)&\geq \sum_u SUNdisponible_u+\sum_u SUNreservado_u.
\end{align}
Las autorizaciones, la unicidad de claims, el límite de emisión y las protecciones contra reentrancia complementan estas igualdades de solvencia. El contrato de SUN permite transferencias fungibles; la trazabilidad de inyección y consumo se conserva en los registros de emisión, asignación y retiro, no en la identidad individual de cada unidad del token.
